\chapter{INTRODUÇÃO}\label{cap:introducao}

Veículos Aéreos Não Tripulados (VANTs) são empregados em inspeção, mapeamento, monitoramento e aplicações logísticas. No contexto de entregas urbanas, o planejamento de rotas e a segurança operacional são questões centrais, pois a missão precisa respeitar restrições de energia, tempo, ambiente e regulamentação \cite{otto2018,rinaldi2024}. No Brasil, operações de aeronaves não tripuladas de uso civil estão sujeitas aos requisitos gerais da Agência Nacional de Aviação Civil (ANAC) e aos procedimentos de acesso ao espaço aéreo definidos pelo Departamento de Controle do Espaço Aéreo (DECEA) \cite{anac2026,decea2026}. Este trabalho não propõe operação comercial ou urbana em espaço aéreo aberto; utiliza o cenário logístico como motivação para o desenvolvimento de uma plataforma experimental em ambiente controlado.

Um quadricóptero autônomo combina problemas de naturezas distintas. O planejamento de missão requer selecionar uma sequência de pontos que reduza o custo de deslocamento, problema relacionado ao Problema do Caixeiro Viajante. A estabilidade da aeronave, por sua vez, exige controle em tempo real de um sistema subatuado, não linear e sujeito a perturbações \cite{beard2012,quan2017}. A integração dessas funções demanda uma arquitetura em que o planejamento não comprometa as funções críticas de estimação, controle e segurança.

Essa integração é relevante porque muitos estudos tratam separadamente o planejamento de rotas, o controle de atitude e a comunicação embarcada. Em um protótipo físico, entretanto, essas partes compartilham restrições de processamento, temporização, sensores, alimentação elétrica e tolerância a falhas. Uma rota calculada corretamente perde valor se não puder ser traduzida em comandos seguros, e um controlador estável em simulação precisa ser confrontado com atrasos, saturações e ruídos presentes no sistema real.

Outra motivação do trabalho é a reprodutibilidade de uma arquitetura experimental de baixo custo relativo. Plataformas abertas, como STM32, Raspberry Pi, FreeRTOS e MAVLink, permitem documentação detalhada das interfaces e dos procedimentos de teste. Essa escolha favorece a análise acadêmica do sistema, pois torna possível registrar versões de hardware, bibliotecas, ferramentas, parâmetros de controle e critérios de avaliação.

O objetivo geral deste trabalho é desenvolver uma plataforma experimental de VANT para investigar a integração entre planejamento de rotas e controle de voo autônomo. A estratégia de controle adotada é o Regulador Linear Quadrático Gaussiano (LQG), que combina um Regulador Linear Quadrático (LQR), projetado a partir da linearização do modelo do quadricóptero em torno da condição de voo estacionário, com um filtro de Kalman para a estimação dos estados. A escolha permite explicitar o compromisso entre desvios de estado e esforço de atuação por meio das matrizes de ponderação do problema de controle ótimo \cite{jafari2010,okyere2019}.

Os objetivos específicos são: modelar a cinemática e a dinâmica do quadricóptero; construir e simular o controlador LQG; implementar a arquitetura embarcada em STM32F411 com FreeRTOS; integrar a comunicação entre o controlador de voo e o Raspberry Pi Zero W; estruturar um Algoritmo Genético para planejamento de rotas; e definir ensaios de bancada, simulação e voo controlado para avaliar o sistema. A validação em voo só foi realizada depois das verificações elétricas, de comunicação e de segurança descritas neste trabalho.

A arquitetura proposta separa a camada de voo da camada de missão. O STM32F411 executa leitura de sensores, estimação, controle LQG, mistura de motores, comando DShot e \textit{failsafes}. O Raspberry Pi Zero W realiza funções de missão, rede, telemetria e registro. A comunicação entre essas camadas é especificada como DroneLink 1.0, protocolo próprio do projeto baseado em UART e enquadramento MAVLink 2 \cite{koubaa2019,mavlinkguide}. O Raspberry Pi envia objetivos de missão e telemetria de alto nível, enquanto o STM32 permanece responsável pelas ações críticas de controle e segurança. Para a estimativa de atitude, o projeto adota um módulo da família BNO08X; o MTF-01 fornece medições de fluxo óptico e distância utilizadas nos testes internos \cite{bno08xdatasheet,micoair2024}.

Como resultado, o trabalho apresenta uma base para experimentos de autonomia em ambiente controlado, com separação clara entre missão e controle de voo. Também foram produzidos artefatos de apoio à continuidade do projeto, incluindo modelo matemático, critérios de sintonia do controlador LQG, especificação de comunicação, configuração dos periféricos e sequência de ensaios. Esses resultados não substituem certificação, validação operacional externa ou atendimento a requisitos regulatórios de voo em espaço aéreo aberto.

O método combina pesquisa bibliográfica, modelagem, simulação, desenvolvimento de \textit{firmware} e experimentação progressiva. As evidências de desempenho são tratadas apenas a partir de simulações, ensaios de bancada e registros experimentais definidos no procedimento metodológico.

Além desta introdução, o Capítulo~\ref{cap_revisao} apresenta os fundamentos de quadricópteros, controle LQG, planejamento de rotas e sistemas embarcados. O Capítulo~\ref{cap:metodologia} descreve os materiais, o ambiente de desenvolvimento, os procedimentos e os critérios de avaliação. O Capítulo~\ref{cap:resultados} apresenta e discute os resultados dos sensores, das simulações, do voo, da comunicação e do planejamento de rotas. O Capítulo~\ref{cap:conclusao} reúne as considerações finais e os trabalhos futuros.
